\section{什么是 Autonomous Agent：从 buzzword 回到定义}

Stone 给出的定义非常朴素：\textbf{agent 是通过 sensors 和 actuators 与环境持续交互的智能系统}。这一定义故意避免把 Agent 等同于聊天窗口里的 API 调用器，也避免把 Agent 仅仅理解为多工具链路。

\begin{knowledgebox}{连续闭环是关键}
在这一讲里，Agent 的最小闭环是：
\[
\text{observe} \rightarrow \text{reason / decide} \rightarrow \text{act} \rightarrow \text{new observation}
\]
重点在 ``continuous loop''，而不是 request-response。
\end{knowledgebox}

按照这一定义，chatbot 在某些设置下可以是 Agent，但很多典型 chatbot 仍偏 ``被动响应''。Autonomous Agent 更强调主动持续行为，不需要每一步都等用户明确触发。

\begin{importantbox}{Robot 与 Agent 的关系}
Stone 用了一个非常实用的区分：
\begin{itemize}
    \item Robot 可以看作是 \textbf{physical agent}；
    \item 所有 robot 都是 agent；
    \item 但并非所有 agent 都必须具备 physical embodiment。
\end{itemize}
这让我们能同时讨论 software agent、game-playing agent、以及 real-world embodied agent。
\end{importantbox}

\begin{warningbox}{常见误解}
把 Agent 缩减为 ``LLM + tool call'' 会直接丢掉两类关键问题：
\begin{itemize}
    \item 真实动作约束（安全、延迟、失败恢复）；
    \item 多体互动约束（协作、竞争、信用分配）。
\end{itemize}
这两类问题恰好决定系统能否从 demo 走向部署。
\end{warningbox}

\subsection{本章小结}
Agent 的定义不是为了学术形式化，而是决定你后续系统边界怎么画。定义太窄，后面的 benchmark 和工程结论都会失真。

